% GENERATED FILE. Edit canonical lessons, then run npm run generate:books.
% canonical-source-hash: fnv1a64:43bda8bc26240aec
% canonical-lessons: SA-C370-sthitih, SA-C370-avastha, SA-C370-parinamati, SA-C370-anukulah, SA-C370-pratikulah

\chapter{As Things Stand}
\label{ch:sa-as-things-stand}
\hlchaptermodality{370}

\begin{chapteropening}
\textbf{By the end of this chapter:} \emph{I can say what I will do if the situation changes.}

Complete the last lesson of chapter 370: I can say what I will do if the situation changes.
\end{chapteropening}

\section[\texorpdfstring{\sk{स्थितिः} (sthitiḥ)}{स्थितिः (sthitiḥ)}]{\sk{स्थितिः} (sthitiḥ) — a situation}

\label{lesson:SA-C370-sthitih}



\begin{quote}
Before the new one: say the Sanskrit for comes by, then the Sanskrit for a likelihood.
\end{quote}

\subsection*{You'll want to know: \sk{स्थितिः}}
\textbf{\sk{स्थितिः}} — \emph{sthitiḥ} — \textquotedblleft{}a situation\textquotedblright{}.

\begin{quote}
\sk{यदि} \sk{स्थितिः} \sk{कठिना} \sk{भवति}, \sk{तर्हि} \sk{अहं} \sk{साहाय्यं} \sk{करिष्यामि।}
\end{quote}

\begin{quote}
\emph{yadi sthitiḥ kaṭhinā bhavati, tarhi ahaṃ sāhāyyaṃ kariṣyāmi.}
\end{quote}

\begin{quote}
If the situation becomes hard, I will help.
\end{quote}

A feminine noun.

\begin{cousinweb}
\textbf{\sk{स्थितिः}} (\emph{sthitiḥ}) comes from the root √sthā, \textquotedblleft{}to stand\textquotedblright{}; MW gives \emph{sthiti} \textquotedblleft{}staying or being in any state or condition\textquotedblright{} (MW).
\end{cousinweb}

\begin{grammarlens}[title={What you've built}]
One more word for saying what happens if.
\end{grammarlens}

\subsection*{Your turn}
\begin{itemize}
\raggedright
  \item \emph{Say it:} \emph{sthitiḥ}
  \item \emph{Say it:} \emph{sthitiḥ}, once more
  \item \emph{Say it:} \emph{yadi sthitiḥ kaṭhinā bhavati, tarhi ahaṃ sāhāyyaṃ kariṣyāmi.} — If the situation becomes hard, I will help.
  \item \emph{Recall:} say the Sanskrit for comes by, then the Sanskrit for a likelihood, then say \emph{sthitiḥ} again
\end{itemize}

\begin{tcolorbox}[breakable,colback=teal!4,colframe=teal!35!black,title={Before you move on}]
What does \sk{स्थितिः} mean? (\textquotedblleft{}A situation\textquotedblright{}.) Say it once more.
\end{tcolorbox}
\section[\texorpdfstring{\sk{अवस्था} (avasthā)}{अवस्था (avasthā)}]{\sk{अवस्था} (avasthā) — a state; a condition}

\label{lesson:SA-C370-avastha}



\begin{quote}
Before the new one: say the Sanskrit for a likelihood, then the Sanskrit for a situation.
\end{quote}

\subsection*{You'll want to know: \sk{अवस्था}}
\textbf{\sk{अवस्था}} — \emph{avasthā} — \textquotedblleft{}a state; a condition\textquotedblright{}.

\begin{quote}
\sk{यदि} \sk{वृद्धस्य} \sk{अवस्था} \sk{उत्तमा}, \sk{तर्हि} \sk{सः} \sk{श्वः} \sk{गमिष्यति।}
\end{quote}

\begin{quote}
\emph{yadi vṛddhasya avasthā uttamā, tarhi saḥ śvaḥ gamiṣyati.}
\end{quote}

\begin{quote}
If the old man's condition is good, he will go tomorrow.
\end{quote}

A feminine noun.

\begin{cousinweb}
\textbf{\sk{अवस्था}} (\emph{avasthā}) is built on \emph{ava}, \textquotedblleft{}down\textquotedblright{}, and the root √sthā, \textquotedblleft{}to stand\textquotedblright{}; MW gives \emph{ava-sthā} \textquotedblleft{}state, condition, situation\textquotedblright{} (MW).
\end{cousinweb}

\begin{grammarlens}[title={What you've built}]
One more word for saying what happens if.
\end{grammarlens}

\subsection*{Your turn}
\begin{itemize}
\raggedright
  \item \emph{Say it:} \emph{avasthā}
  \item \emph{Say it:} \emph{avasthā}, once more
  \item \emph{Say it:} \emph{yadi vṛddhasya avasthā uttamā, tarhi saḥ śvaḥ gamiṣyati.} — If the old man's condition is good, he will go tomorrow.
  \item \emph{Recall:} say the Sanskrit for a likelihood, then the Sanskrit for a situation, then say \emph{avasthā} again
\end{itemize}

\begin{tcolorbox}[breakable,colback=teal!4,colframe=teal!35!black,title={Before you move on}]
What does \sk{अवस्था} mean? (\textquotedblleft{}A state; a condition\textquotedblright{}.) Say it once more.
\end{tcolorbox}
\section[\texorpdfstring{\sk{परिणमति} (pariṇamati)}{परिणमति (pariṇamati)}]{\sk{परिणमति} (pariṇamati) — turns into; changes into}

\label{lesson:SA-C370-parinamati}



\begin{quote}
Before the new one: say the Sanskrit for a situation, then the Sanskrit for a state of things.
\end{quote}

\subsection*{You'll want to know: \sk{परिणमति}}
\textbf{\sk{परिणमति}} — \emph{pariṇamati} — \textquotedblleft{}turns into; changes into\textquotedblright{}.

\begin{quote}
\sk{यदि} \sk{शीतं} \sk{भवति}, \sk{तर्हि} \sk{जलं} \sk{हिमेन} \sk{परिणमति।}
\end{quote}

\begin{quote}
\emph{yadi śītaṃ bhavati, tarhi jalaṃ himena pariṇamati.}
\end{quote}

\begin{quote}
If it is cold, water turns into ice.
\end{quote}

Past (it …): \sk{जलं} \sk{हिमेन} \sk{परिणतवत्}, \emph{jalaṃ himena pariṇatavat}.

\begin{cousinweb}
\textbf{\sk{परिणमति}} (\emph{pariṇamati}) is built on \emph{pari}, \textquotedblleft{}around\textquotedblright{}, and the root √nam, \textquotedblleft{}to bend\textquotedblright{}; MW gives \emph{pari-√ṇam} \textquotedblleft{}to bend aside\textquotedblright{} and \textquotedblleft{}to change or be transformed into\textquotedblright{} (MW).
\end{cousinweb}

\begin{grammarlens}[title={What you've built}]
One more word for saying what happens if.
\end{grammarlens}

\subsection*{Your turn}
\begin{itemize}
\raggedright
  \item \emph{Say it:} \emph{pariṇamati}
  \item \emph{Say it:} \emph{pariṇamati}, once more
  \item \emph{Say it:} \emph{yadi śītaṃ bhavati, tarhi jalaṃ himena pariṇamati.} — If it is cold, water turns into ice.
  \item \emph{Recall:} say the Sanskrit for a situation, then the Sanskrit for a state of things, then say \emph{pariṇamati} again
\end{itemize}

\begin{tcolorbox}[breakable,colback=teal!4,colframe=teal!35!black,title={Before you move on}]
What does \sk{परिणमति} mean? (\textquotedblleft{}Turns into; changes into\textquotedblright{}.) Say it once more.
\end{tcolorbox}
\section[\texorpdfstring{\sk{अनुकूलः} (anukūlaḥ)}{अनुकूलः (anukūlaḥ)}]{\sk{अनुकूलः} (anukūlaḥ) — favourable}

\label{lesson:SA-C370-anukulah}



\begin{quote}
Before the new one: say the Sanskrit for a state of things, then the Sanskrit for turns into.
\end{quote}

\subsection*{You'll want to know: \sk{अनुकूलः}}
\textbf{\sk{अनुकूलः}} — \emph{anukūlaḥ} — \textquotedblleft{}favourable\textquotedblright{}.

\begin{quote}
\sk{यदि} \sk{वायुः} \sk{अनुकूलः}, \sk{तर्हि} \sk{नौका} \sk{शीघ्रं} \sk{गमिष्यति।}
\end{quote}

\begin{quote}
\emph{yadi vāyuḥ anukūlaḥ, tarhi naukā śīghraṃ gamiṣyati.}
\end{quote}

\begin{quote}
If the wind is favourable, the boat will go fast.
\end{quote}

\begin{cousinweb}
\textbf{\sk{अनुकूलः}} (\emph{anukūlaḥ}) is built from \emph{anu}, \textquotedblleft{}along\textquotedblright{}, and \emph{kūla}, \textquotedblleft{}a slope, a bank\textquotedblright{}; MW gives \emph{anu-kūla} \textquotedblleft{}favourable, agreeable\textquotedblright{} (MW).
\end{cousinweb}

\begin{grammarlens}[title={What you've built}]
One more word for saying what happens if.
\end{grammarlens}

\subsection*{Your turn}
\begin{itemize}
\raggedright
  \item \emph{Say it:} \emph{anukūlaḥ}
  \item \emph{Say it:} \emph{anukūlaḥ}, once more
  \item \emph{Say it:} \emph{yadi vāyuḥ anukūlaḥ, tarhi naukā śīghraṃ gamiṣyati.} — If the wind is favourable, the boat will go fast.
  \item \emph{Recall:} say the Sanskrit for a state of things, then the Sanskrit for turns into, then say \emph{anukūlaḥ} again
\end{itemize}

\begin{tcolorbox}[breakable,colback=teal!4,colframe=teal!35!black,title={Before you move on}]
What does \sk{अनुकूलः} mean? (\textquotedblleft{}Favourable\textquotedblright{}.) Say it once more.
\end{tcolorbox}
\section[\texorpdfstring{\sk{प्रतिकूलः} (pratikūlaḥ)}{प्रतिकूलः (pratikūlaḥ)}]{\sk{प्रतिकूलः} (pratikūlaḥ) — adverse; unfavourable}

\label{lesson:SA-C370-pratikulah}



\begin{quote}
Before the new one: say the Sanskrit for turns into, then the Sanskrit for favourable.
\end{quote}

\subsection*{You'll want to know: \sk{प्रतिकूलः}}
\textbf{\sk{प्रतिकूलः}} — \emph{pratikūlaḥ} — \textquotedblleft{}adverse; unfavourable\textquotedblright{}.

\begin{quote}
\sk{यदि} \sk{वायुः} \sk{प्रतिकूलः}, \sk{तर्हि} \sk{नौका} \sk{न} \sk{गमिष्यति।}
\end{quote}

\begin{quote}
\emph{yadi vāyuḥ pratikūlaḥ, tarhi naukā na gamiṣyati.}
\end{quote}

\begin{quote}
If the wind is against us, the boat will not go.
\end{quote}

\begin{cousinweb}
\textbf{\sk{प्रतिकूलः}} (\emph{pratikūlaḥ}) is built from \emph{prati}, \textquotedblleft{}against\textquotedblright{}, and \emph{kūla}, \textquotedblleft{}a bank\textquotedblright{}; MW glosses it as \textquotedblleft{}against the bank\textquotedblright{} and gives \textquotedblleft{}contrary, adverse, opposite\textquotedblright{} (MW).
\end{cousinweb}

\begin{grammarlens}[title={What you've built}]
One more word for saying what happens if.
\end{grammarlens}

\subsection*{Your turn}
\begin{itemize}
\raggedright
  \item \emph{Say it:} \emph{pratikūlaḥ}
  \item \emph{Say it:} \emph{pratikūlaḥ}, once more
  \item \emph{Say it:} \emph{yadi vāyuḥ pratikūlaḥ, tarhi naukā na gamiṣyati.} — If the wind is against us, the boat will not go.
  \item \emph{Recall:} say the Sanskrit for turns into, then the Sanskrit for favourable, then say \emph{pratikūlaḥ} again
\end{itemize}

\begin{tcolorbox}[breakable,colback=teal!4,colframe=teal!35!black,title={Before you move on}]
What does \sk{प्रतिकूलः} mean? (\textquotedblleft{}Adverse; unfavourable\textquotedblright{}.) Say it once more.
\end{tcolorbox}
